\section{Accent Robustness in Speech Translation}
\label{sec:accent-st}

Section III established that accent-related degradation in automatic speech recognition (ASR) is systematic, measurable, and present across the open model families relevant to this project, including Whisper and SeamlessM4T. That evidence, however, is confined to ASR-level metrics such as word error rate (WER) and word information lost (WIL). It does not by itself establish whether accent-driven recognition error affects final translation quality, or whether the two architectural paradigms compared in this survey are affected to a similar degree. This section reviews work at the intersection of accent, error propagation, and speech translation (ST) to test whether that specific combination has already been addressed.

One important area of research concerns the effect of ASR errors on cascaded ST. Sperber \textit{et al.} \cite{sperber2019} investigated this by directly manipulating the intermediate ASR transcripts provided to different translation systems. They compared systems using correct reference transcripts, normal ASR output, and transcripts containing artificially introduced substitution, insertion, and deletion noise. They report that the cascade responds most strongly to both improved and degraded labels, confirming that it is the architecture most exposed to intermediate transcription error, while models with continuous access to acoustic information are comparatively, though not completely, insulated.

Nguyen \textit{et al.} \cite{pida2026} quantify this effect directly for Vietnamese-English cascaded ST, reporting BLEU drops of 6.79 to 10.64 points when translating ASR output rather than clean reference transcripts. Using linear mixed-effects modelling, they further show that, among substitution-type errors specifically, most arise from a systematic phonetic cause involving vowels, consonants, and tones rather than from phonetically unrelated substitutions. The phonetic-confusion categories are statistically significant predictors of translation-quality degradation, whereas phonetically unrelated substitutions are not. Insertions and deletions, however, remain the largest-magnitude predictors of translation-quality degradation overall. Read together with the AESRC2020 and accent-codebook evidence from Section III, this is a relevant methodological bridge: if accent variation is itself a phonetically systematic source of ASR error, and phonetically systematic substitution errors are significant (if not the largest-magnitude) contributors to cascaded translation-quality loss, then the mechanism by which accent could degrade cascaded ST specifically is plausible and consistent with existing evidence, even though neither study manipulates accent as an independent variable.

A second group of studies examines the robustness of speech models to other forms of variation. The SeamlessM4T technical report \cite{seamlessm4t2023} and its successor \cite{seamless2023} benchmark robustness to background noise and to what the authors label ``speaker variations'' or ``vocal style variations,'' reporting that SeamlessM4T outperforms a Whisper-Large-v2 baseline by a substantial margin on both axes, using BLEU/chrF-based metrics and a coefficient-of-variation measure across speaking styles. This is the closest published evidence to a controlled Direct-versus-Whisper robustness comparison among the model families used in this project. Two limitations constrain how far it can be read as evidence for the present research question. First, ``speaker variation'' in these reports is not equivalent to, and is not reported as, regional accent; the test conditions are not accent-labelled, so any accent effect is at best a subset of an unspecified broader category. Second, the comparison evaluates Whisper alone rather than a full Whisper-to-MT cascade of the kind this project benchmarks, so a cascade-specific error-propagation effect of the kind documented by Sperber \textit{et al.} is not isolated in these results. The reported robustness advantage of SeamlessM4T over Whisper therefore cannot be attributed to accent specifically, nor to the direct-versus-cascade distinction as operationalised in this survey.

A narrower but genuinely relevant example comes from Havare \textit{et al.} \cite{losttranscription2026}, who study a multilingual, code-mixed speech-driven framework for programming queries. In their qualitative failure-mode taxonomy, the authors explicitly name ``accent based errors'' as one of seven ASR failure modes, giving concrete regional examples: ``default'' pronounced with a Marathi accent is transcribed as ``de fall,'' and ``machine'' pronounced with a Telugu accent becomes ``mishan.'' This is the clearest example identified in this survey of accent named explicitly, with concrete regional examples, as a cause of ASR misrecognition within a multi-stage pipeline. The paper also quantifies downstream propagation directly: on a code-retrieval task, Recall@5 drops from 92\% using gold-text queries to 87\% using raw ASR-transcribed queries, and MRR drops from 87.2\% to 81.58\%, before an LLM-based correction stage partially recovers performance. Two limitations nonetheless constrain how far this can be read as evidence for the present research question. First, the accent-based failure mode is presented qualitatively, with two illustrative examples, rather than as a controlled variable in the paper's quantitative evaluation. The reported WER, PER, WFED, and Recall/MRR results are instead broken down by target language, not by accent group. Therefore, the quantified downstream-degradation numbers cannot themselves be attributed to accent specifically. Second, the domain is code-mixed spoken programming queries processed through a Whisper/Indic-Conformer ASR stage feeding an LLM-based retrieval pipeline, not natural spoken-English translation compared across a Direct-versus-Cascaded ST architecture; no direct-style architecture is evaluated as a comparator in the sense used in this survey. The paper is therefore genuine, if partial, evidence that accent-driven ASR error can propagate into degraded downstream task performance in at least one adjacent pipeline, but it does not establish this for speech translation specifically, nor for a Direct-versus-Cascaded comparison.

Read together, these strands of literature support two provisional conclusions. First, the mechanism through which accent-related ASR errors could affect cascaded translation quality is consistent with existing evidence. Error propagation in cascaded ST has been directly demonstrated, and phonetically systematic substitution errors have been shown to contribute significantly to translation-quality degradation. The evidence reviewed in Section III further establishes accent variation as a potential source of systematic phonetic differences, while Havare \textit{et al.} \cite{losttranscription2026} explicitly identify accent-related ASR errors and demonstrate their propagation into measurable downstream task degradation in an adjacent speech-processing application. The existing literature therefore provides a plausible theoretical and empirical basis for investigating accent-related degradation in cascaded ST. Second, no study identified in this survey evaluates matched Direct and Cascaded ST systems on the same accent-labelled regional English speech, using final translation quality as the outcome, with accent treated as a controlled rather than illustrative variable. The closest available architecture-level robustness comparison \cite{seamlessm4t2023,seamless2023} does not use accent-labelled data or a matched cascade condition, and the closest available accent-explicit propagation evidence \cite{losttranscription2026} names accent only qualitatively and outside the speech-translation domain. The combination of these elements, namely accent as a controlled variable, a matched Direct/Cascaded ST comparison, and final translation quality as the outcome, therefore represents a potential research gap. This gap must, however, be considered alongside the remaining themes reviewed in Section V before it can be established as the central research gap of the project.
